% Part: methods
% Chapter: induction
% Section: relations

\documentclass[../../../include/open-logic-section]{subfiles}

\begin{document}

\olfileid{mth}{ind}{rel}

\olsection{સંબંધો અને વિધેયો}

પદાર્થોનો ગણ (જેમ પ્રાકૃતિક સંખ્યાઓ અથવા સુઘડ પદો)
આગમનાત્મક રીતે વ્યાખ્યાયિત કર્યો હોય ત્યારે આ પદાર્થો
\emph{પરના સંબંધો} પણ આગમનથી વ્યાખ્યાયિત કરી શકીએ.
દાખલા તરીકે, આ વિચાર લો: સુઘડ પદ~$t_1$ જો સુઘડ
પદ~$t_2$ના ભાગરૂપે આવતું હોય તો તે તેનું ઉપપદ છે.
તે માટે $t_1 \sqsubseteq t_2$ ચિહ્ન વાપરીએ.
અલબત્ત, દરેક સુઘડ પદ પોતાનું ઉપપદ છે:
$t \sqsubseteq t$. આ સંબંધની આગમનાત્મક વ્યાખ્યા
આ રીતે આપી શકીએ:

\begin{defn}
  સુઘડ પદ~$t_1$ એ~$t_2$નું ઉપપદ છે એવો સંબંધ,
  $t_1 \sqsubseteq t_2$,~$t_2$ પરના આગમનથી
  નીચે મુજબ વ્યાખ્યાયિત થાય છે:
  \begin{enumerate}
  \item $t_2$ અક્ષર હોય તો $t_1 \sqsubseteq t_2$
    ત્યારે અને ત્યારે જ જ્યારે $t_1 = t_2$.
  \item $t_2$ એ $[s_1 \circ s_2]$ હોય તો
    $t_1 \sqsubseteq t_2$ ત્યારે અને ત્યારે જ જ્યારે
    $t_1 = t_2$, $t_1 \sqsubseteq s_1$ અથવા
    $t_1 \sqsubseteq s_2$.
  \end{enumerate}
\end{defn}

આ વ્યાખ્યાથી, દાખલા તરીકે, $\mathrm{a}
\sqsubseteq [\mathrm{b} \circ \mathrm{a}]$
મળે છે. કારણ કે (2) કહે છે કે
$\mathrm{a} \sqsubseteq [\mathrm{b} \circ \mathrm{a}]$
ત્યારે અને ત્યારે જ જ્યારે
$\mathrm{a} = [\mathrm{b} \circ \mathrm{a}]$,
અથવા $\mathrm{a} \sqsubseteq \mathrm{b}$,
અથવા $\mathrm{a} \sqsubseteq \mathrm{a}$.
પહેલી બે વાતો ખોટી છે: $\mathrm{a}$ એ
$[\mathrm{b} \circ
  \mathrm{a}]$ નથી, અને (1) મુજબ
$\mathrm{a} \sqsubseteq \mathrm{b}$
ત્યારે અને ત્યારે જ જ્યારે $\mathrm{a} = \mathrm{b}$,
જે પણ ખોટું છે. પરંતુ (1) મુજબ જ
$\mathrm{a} \sqsubseteq \mathrm{a}$
ત્યારે અને ત્યારે જ જ્યારે $\mathrm{a} = \mathrm{a}$,
જે સાચું છે.

ધ્યાન આપવાની અગત્યની વાત એ છે કે આ વ્યાખ્યા
સફળ થવી એક એવી હકીકત પર આધાર રાખે છે જે
આપણે હજી સાબિત કરી નથી: દરેક સુઘડ પદ~$t$
કાં તો પોતે એક અક્ષર હોય છે, અથવા
\emph{એકમાત્ર રીતે નિર્ધારિત} સુઘડ પદો
$s_1$ અને $s_2$ મળે છે જેથી
$t = [s_1 \circ s_2]$. અહીં "એકમાત્ર રીતે
નિર્ધારિત"નો અર્થ એ છે કે
$t = [s_1 \circ s_2]$ હોય તો
$s_1 \neq r_1$ અથવા $s_2 \neq r_2$ હોય
એવા બીજા બે પદો માટે તે \emph{સાથે સાથે}
$= [r_1 \circ r_2]$ પણ નહીં હોય. જો એવું
હોય તો ખંડ~(2) પોતાનાં જ પરિણામો સાથે
વિરોધમાં આવી શકે:~$t_2$ને
$[s_1 \circ s_2]$ તરીકે વાંચીએ ત્યારે
$t_1 \sqsubseteq t_2$ મળે, પણ~$t_2$ને
$[r_1 \circ r_2]$ તરીકે વાંચીએ ત્યારે
$t_1 \sqsubseteq t_2$ ન મળે. આવું ન
થઈ શકે તે સાબિત કરતાં પહેલાં જ્યાં
આવું \emph{થઈ શકે} તેવું ઉદાહરણ જોઈએ.

\begin{defn}
  \emph{કૌંસવિહોણાં પદો} આગમનાત્મક રીતે
  આ પ્રમાણે વ્યાખ્યાયિત કરો:
  \begin{enumerate}
  \item દરેક અક્ષર કૌંસવિહોણું પદ છે.
    \item $s_1$ અને $s_2$ કૌંસવિહોણાં પદો
      હોય તો $s_1 \circ s_2$ પણ કૌંસવિહોણું
      પદ છે.
    \item આ સિવાય બીજું કંઈ કૌંસવિહોણું પદ નથી.
  \end{enumerate}
\end{defn}

કૌંસવિહોણાં પદોનાં ઉદાહરણો
$\mathrm{a}$, $\mathrm{b} \circ
\mathrm{d}$ અને $\mathrm{b} \circ \mathrm{a}
\circ \mathrm{b}$ છે. હવે કૌંસવિહોણાં
પદો માટે પણ ઉપરની જેમ "ઉપપદ" વ્યાખ્યાયિત
કરીએ તો બીજો ખંડ આવો થશે:
\begin{center}
  જો $t_2 = s_1 \circ s_2$ હોય, તો
  $t_1 \sqsubseteq t_2$ ત્યારે અને ત્યારે જ
  જ્યારે $t_1 = t_2$, $t_1
  \sqsubseteq s_1$ અથવા $t_1 \sqsubseteq s_2$.
\end{center}

હવે $\mathrm{b} \circ \mathrm{a} \circ \mathrm{b}$
એ $s_1 \circ s_2$ સ્વરૂપનું છે, જ્યાં
\begin{align*}
  s_1 & = \mathrm{b} \text{ અને} & s_2 & = \mathrm{a} \circ \mathrm{b}.
\intertext{તે $r_1 \circ r_2$ સ્વરૂપનું પણ છે, જ્યાં}
  r_1 & = \mathrm{b} \circ \mathrm{a} \text{ અને} & r_2 &= \mathrm{b}.
\end{align*}
હવે $\mathrm{a} \circ \mathrm{b}$ એ
$\mathrm{b} \circ
\mathrm{a} \circ \mathrm{b}$નું ઉપપદ છે?
પહેલા વાંચન પ્રમાણે જવાબ હા છે, અને બીજા
વાંચન પ્રમાણે ના છે.

સુઘડ પદને બીજાં સુઘડ પદોમાંથી બનાવવાની રીત
એકમાત્ર હોય તે ગુણધર્મને \emph{એકમાત્ર વાંચન}
કહે છે. આવા આગમનાત્મક રીતે વ્યાખ્યાયિત
પદાર્થો પર સંબંધોની આગમનાત્મક વ્યાખ્યાઓ
મહત્વની હોવાથી આ ગુણધર્મ સાબિત કરવો પડે.

\begin{prop}
  ધારો કે $t$ સુઘડ પદ છે. ત્યારે કાં તો $t$
  પોતે એક અક્ષર છે, અથવા એકમાત્ર રીતે
  નિર્ધારિત સુઘડ પદો $s_1$, $s_2$ મળે છે
  જેથી~$t = [s_1 \circ s_2]$.
\end{prop}

\begin{proof}
  $t$ પોતે અક્ષર હોય તો શરત પૂરી થાય છે.
  તેથી ધારો કે $t$ પોતે અક્ષર નથી.
  આગમનાત્મક વ્યાખ્યા મુજબ ત્યારે કોઈ
  સુઘડ પદો $s_1$ અને~$s_2$ માટે $t$
  એ $[s_1 \circ s_2]$ સ્વરૂપનું હોવું જોઈએ.
  હવે આ બંને એકમાત્ર રીતે નિર્ધારિત છે
  તે બતાવવું બાકી છે: જો
  $t = [r_1 \circ r_2]$ હોય તો
  $s_1 = r_1$ અને $s_2 = r_2$.

  તેથી ધારો કે સુઘડ પદો $s_1$, $s_2$,
  $r_1$, $r_2$ માટે
  $t = [s_1 \circ s_2]$ અને
  $t = [r_1 \circ r_2]$ બંને છે.
  $s_1 = r_1$ અને $s_2 = r_2$ બતાવવું
  છે. પહેલા $s_1$ અને $r_1$ એકસમાન
  હોવા જોઈએ, કારણ કે નહિ તો એક બીજાનો
  ઉચિત આરંભિક ખંડ છે. પરંતુ
  \olref[sti]{prop:initial} મુજબ $s_1$
  અને~$r_1$ બંને સુઘડ પદો હોય ત્યારે
  એવું અશક્ય છે.
  હવે $s_1 = r_1$ હોય તો સ્પષ્ટપણે
  $s_2 = r_2$ પણ છે.
\end{proof}

વિધેયો પણ આગમનાત્મક રીતે વ્યાખ્યાયિત કરી
શકીએ. દાખલા તરીકે, કોઈ પણ સુઘડ પદમાં
એકની અંદર એક આવેલા $[\dots]$ કૌંસોનું
મહત્તમ ઊંડાણ આપતું વિધેય~$f$ આ રીતે
વ્યાખ્યાયિત કરીએ:

\begin{defn}
  \ollabel{defn:depth} સુઘડ પદનું
  \emph{ઊંડાણ}, $f(t)$, આગમનાત્મક રીતે
  નીચે મુજબ વ્યાખ્યાયિત થાય છે:
  \[
  f(t) = \begin{cases}
    0 & \text{ જો $t$ અક્ષર હોય}\\
    \max(f(s_1), f(s_2)) + 1 & \text{ જો $t = [s_1 \circ s_2]$ હોય.}
  \end{cases}
  \]
\end{defn}

દાખલા તરીકે,
\begin{align*}
  f([\mathrm{a} \circ \mathrm{b}]) & =
    \max(f(\mathrm{a}),f(\mathrm{b})) + 1 = \\
  &= \max(0, 0) + 1 = 1, \text{ અને}\\
  f([[\mathrm{a} \circ \mathrm{b}] \circ \mathrm{c}]) & =
    \max(f([\mathrm{a} \circ \mathrm{b}]), f(\mathrm{c})) + 1 = \\
  & = \max(1,0) + 1 = 2.
\end{align*}

અહીં આપણે અલબત્ત $s_1$ અને $s_2$ સુઘડ
પદો છે એમ ધારીએ છીએ અને દરેક સુઘડ પદ
કાં તો અક્ષર છે અથવા $[s_1 \circ s_2]$
સ્વરૂપનું છે તે હકીકત વાપરીએ છીએ. આ
સ્વરૂપની રચના એક જ રીતે થઈ શકે એ
અહીં પણ અગત્યનું છે. તેનું કારણ જોવા
અગાઉ વ્યાખ્યાયિત કૌંસવિહોણાં પદો ફરી
વિચારો. તેમને માટે અનુરૂપ "વ્યાખ્યા"
આવી થશે:
\[
  g(t) =
  \begin{cases}
    0 & \text{ જો $t$ અક્ષર હોય}\\
   \max(g(s_1), g(s_2)) + 1 & \text{ જો $t = s_1 \circ s_2$ હોય.}
  \end{cases}
\]
હવે કૌંસવિહોણું પદ
$\mathrm{a} \circ \mathrm{b} \circ
\mathrm{c} \circ \mathrm{d}$ વિચારો.
તેને એક કરતાં વધુ રીતે વાંચી શકાય છે;
દાખલા તરીકે, $s_1 \circ s_2$ સ્વરૂપમાં
\begin{align*}
  s_1 & = \mathrm{a} \text{ અને} &
  s_2 & = \mathrm{b} \circ \mathrm{c} \circ \mathrm{d},
\intertext{અથવા $r_1 \circ r_2$ સ્વરૂપમાં}
  r_1 & = \mathrm{a} \circ \mathrm{b} \text{ અને} &
  r_2 &= \mathrm{c} \circ \mathrm{d}.
\end{align*}
પહેલા વાંચન મુજબ $g$ ગણીએ તો મળે:
\begin{align*}
  g(s_1 \circ s_2) & =
  \max(g(\mathrm{a}), g(\mathrm{b} \circ \mathrm{c} \circ \mathrm{d})) + 1 =\\ & =
  \max(0,2) + 1 = 3
  \intertext{જ્યારે બીજા વાંચન મુજબ મળે}
  g(r_1 \circ r_2) & =
  \max(g(\mathrm{a} \circ \mathrm{b}), g(\mathrm{c} \circ \mathrm{d})) + 1=\\ &
  = \max(1,1) + 1 = 2
\end{align*}
પરંતુ વિધેય હંમેશાં એકમાત્ર મૂલ્ય
આપતું હોવું જોઈએ; તેથી~$g$ની આપણી
"વ્યાખ્યા" ખરેખર વિધેય વ્યાખ્યાયિત જ કરતી નથી.
\sourcecorrection{OLMD-166}{સ્થિર અંગ્રેજી મૂળમાં બે
  જગ્યાએ અક્ષર~$\mathrm{b}$ને ચલ~$b$ની
  જેમ ટાઇપસેટ કર્યું છે: ઉપપદના ઉદાહરણમાં
  અને બીજા કૌંસવિહોણા વાંચનના $r_1$માં.
  મૂળનો અભિપ્રાય અક્ષર~$\mathrm{b}$ જ છે;
  બંને જગ્યાએ તે ચિહ્ન સ્પષ્ટ કર્યું છે.}

\begin{prob}
  વિધેય $l$ની આગમનાત્મક વ્યાખ્યા આપો, જ્યાં
  $l(t)$ એ સુઘડ પદ~$t$માં ચિહ્નોની સંખ્યા છે.
\end{prob}

\begin{prob}
  સુઘડ પદો $t$ પર સંરચનાત્મક આગમનથી
  $f(t) < l(t)$ સાબિત કરો (અહીં $l(t)$
  એ~$t$માં ચિહ્નોની સંખ્યા છે અને $f(t)$
  એ \olref[mth][ind][rel]{defn:depth}માં
  વ્યાખ્યાયિત~$t$નું ઊંડાણ છે).
\end{prob}

\end{document}
